\chapter{级联MUX——4大变式详解}

\section{变式1：用2选1 MUX构建8选1 MUX}

\begin{detailbox}[完整教学]
\textbf{【树形结构】}3层：第一层4个→第二层2个→第三层1个=7个2选1 MUX。

\textbf{【选择位分配】}
\begin{itemize}
  \item S0→第一层(4个MUX)：每两组数据中选一个
  \item S1→第二层(2个MUX)：第一层结果中两两选择
  \item S2→第三层(1个MUX)：第二层两个结果中最终选择
\end{itemize}

\textbf{【通用公式】}$2^n$选1 MUX需要$2^n-1$个2选1 MUX，分$n$层。

\textbf{【VHDL——generate结构级】}
\begin{lstlisting}
-- 递归描述流水线MUX树
signal layer0 : std_logic_vector(3 downto 0);
signal layer1 : std_logic_vector(1 downto 0);
begin
  gen0: for i in 0 to 3 generate
    layer0(i) <= D(2*i) when S(0)='0' else D(2*i+1);
  end generate;
  gen1: for i in 0 to 1 generate
    layer1(i) <= layer0(2*i) when S(1)='0' else layer0(2*i+1);
  end generate;
  Y <= layer1(0) when S(2)='0' else layer1(1);
end behavioral;
\end{lstlisting}
\end{detailbox}


\section{变式2：用4选1 MUX构建16选1 MUX}

\begin{detailbox}[完整教学]
\textbf{【结构】}第一层4片4选1 MUX，第二层1片4选1 MUX。共5片。

S1S0→第一层(4片4选1 MUX)。S3S2→第二层(1片4选1 MUX选择第一层输出之一)。
\end{detailbox}


\section{变式3：不对称级联（构建6选1 MUX）}

\begin{detailbox}[完整教学]
\textbf{【问题】}需要6选1，但MUX只有$2^n$选1。用8选1实现6选1（两个输入悬空或接地）。

\textbf{【VHDL——4选1+2选1】}
\begin{lstlisting}
-- 第一层：4选1处理D0-D3
mux4: entity work.mux4to1 port map(D=>D(3 downto 0), S=>S(1 downto 0), Y=>mux4_out);
-- 顶层：2选1在mux4_out和D4/D5间选择
Y <= mux4_out when S(2)='0' else
     D(4) when S(1 downto 0)="00" else D(5);
-- 注意：S2=1时S1S0只需表示2种状态即可
\end{lstlisting}

\textbf{【简化方案】}直接用8选1 MUX，D6和D7接地(固定为0)。
\end{detailbox}


\section{变式4：MUX树实现逻辑函数}

\begin{detailbox}[完整教学]
\textbf{【本质】}MUX=通用逻辑实现器。MUX树=多层选择=分而治之。

\textbf{【Shannon展开定理】}$F(x_1,x_2,\dots,x_n)= \overline{x_1}\cdot F(0,x_2,\dots,x_n) + x_1\cdot F(1,x_2,\dots,x_n)$

这正是2选1 MUX的表达式！递归应用=构建MUX树。

\textbf{【结论】}任何组合逻辑函数都可用2选1 MUX树实现——这其实就是FPGA中LUT的工作原理。
\end{detailbox}
